\section{What the experiments now explain together}
\label{sec:synthesis}
The initial trace followed a witness through an analytic argument. The encounter with classical choice raised an interpretive question: if the actual retained set cannot be extracted, which parts of the subsequent estimate follow for every allowed witness? The experiments have gradually reduced that question to one that can be answered, in part, by hand. The selection score enforces local counts and marked cells. The gain also depends on the presence of adjacent pairs for specific generators. Local count control and control of a difference set are distinct obligations. The same retained set can satisfy the first and leave a positive proportion of odd allowance recoverable under the second.

\subsection{The sequence of conclusions}
The first positive example guaranteed \(2/45\) with one selector-compatible family. A 32-cycle deletion rule subsequently attained \(10367/124242\) on its worst labeling, or about \(8.3442\%\). It became clear that the ranks cheapest for exposure and the residues best for suppressing edges need not coincide. The previous 64-cycle rule balanced them by keeping a dense quarter for supply and a sparse half for blocking. It passed ten percent by a small explicit surplus on its most difficult decision. The present study identifies two independent sources of slack in that construction. Within a support, rates can be exchanged until a leaf mark and a sparse-parent mark bind. Across supports, exchanging entire residues can erase more directions, often raising the attainable return past eleven percent. \href{https://github.com/nousentllc/riemann-hypothesis/pull/12}{First positive recovery}; \href{https://github.com/nousentllc/riemann-hypothesis/pull/13}{The 32-cycle and the multiscale construction}; \href{https://github.com/nousentllc/riemann-hypothesis/pull/14}{The first 64-cycle certificate}.

These comparisons also establish a limit to the intuition that the largest occupied subset must yield the strongest fallback. An added residue can improve the tile probability by providing an even mark while simultaneously closing the last missing difference that earned a whole subtree its recovery. Adjusting a positive rate changes only the marginal price until the support changes. Removing the last point in a difference class changes the coefficients of several other rates at once. Thus continuous optimization is appropriate inside each support, with a discrete comparison at the boundary. The exact optimum on the old support measures how much one could ever hope to obtain by the continuous part alone.

\subsection{From a percentage back to a gap}
It helps to put the numerator and denominator back in their original units. At the more simply rounded refined profile \(c=2.392096615,\ b=0.706267\), on the literal old adverse labeling we have
\[
     P=277.48320734,\qquad Q=2742.19308421 .
\]
In the normalization derived from the prime ledger, \(Q/(4\cdot64^2)=0.16737018\ldots\) is the limiting odd charge, and \(P/(4\cdot64^2)=0.01693623\ldots\) is the recovered part. The remaining signed charge in that comparison is approximately \(0.15043395\). The ratio of the first two nontrivial quantities is the advertised recovery percentage. It is a property of that particular labeling and constructed family. The all-label theorem guarantees the fraction, while the allowance itself can vary across labelings.

On the same host, an actual selected witness accepting all the source obligations need not realize this arithmetic ledger. The written argument produces another witness, uses it in both gains, and controls what was previously an unquantified degree of freedom. This is a concrete answer to an inference about the proof: a score small enough for selection does not, by itself, force the fallback to be negligible relative to the surviving odd charge. It is not an improvement to the theorem about zeros. The measured percentages concern the accounting of witnesses compatible with one step in a proof.

\subsection{Why a fixed resolution has a ceiling}
A positive-rate residue ultimately supplies adjacent pairs on its fine-scale allowed directions in the sampled construction. A fixed quotient can therefore block edges only on its finitely many coarser splits. The baseline tail in \(Q\) survives even when all available coarse directions disappear. A window strong enough to permit the selection also demands even mass. The optimizer pays for that mass in two ways. It spends some of its fixed total on leaves with little direct benefit, and the physical presence of those leaves can open differences that eliminate benefit elsewhere. The adverse-label upper certificate makes both mechanisms comparable to a concrete lower example.

This explanation is compatible with the preceding multiscale result, where patterns at increasingly many scales can make the recovered share arbitrarily close to one. There the optimizer changes the spatial resolution as the family changes and can suppress allowed pairs that a homogeneous 64-cell model inevitably realizes. Fixing a resolution before taking the arithmetic limit is a mathematical hypothesis. A ceiling for that fixed resolution neither invalidates the multiscale construction nor estimates the optimum of correlated sampling laws.

\subsection{Which part of the research question is now closed?}
The full-support computation collapsed the minimax question at fixed resolution to two masks and then to two binding rates. The open numerical interval in the earlier, globally priced upper certificate reflected a relaxation: one set of leaf prices had to control incompatible geometries at once. Once the structural support was fixed, its expensive singleton and its alternative count-adjustment leaf could be valued at their actual coefficients. The resulting upper equals the constructive lower. Improving this robust value requires changing at least one of the hypotheses.

There are still two natural mathematical questions. The first is whether a retained family can pass the original score while systematically missing the literal product-mark threshold used here to prove feasibility. A necessary estimate directly from the score, strong enough to recover or disprove the limiting cap and the singleton requirement, would widen the theorem. The second is whether allowed edges can be kept rare through correlated choices within a single residue. Precisely where the homogeneous exact-count law made an allowed step realize many adjacent pairs, a correlated law would have to suppress them without causing the marked windows to fail. Either experiment has to keep its witness fixed across all primes and lifts. A better fixed-cycle value in the wider class would identify which obligation in the present theorem is doing real work.

There is also a lower-cost question on the formal side. The final finite result has an unusually small interface: one adverse word, an integer overflow bound, a complete exhaustion down to two masks, and one inequality for each. Replacing the two independent executable enumerations by a formally verified checker consuming a compact certificate would narrow the verification gap. This differs from a full re-formalization of the analytic transport. It specifies a concrete next theorem and exactly which current claims its proof would strengthen.

There is a methodological connection to prior work on revisable memory. In both settings a representation sufficient to score the present output may discard information needed for the next intervention. Here two children with identical occupancy counts can have different sets of viable successors. The explicit example in Section~\ref{sec:geometry} demonstrates the failure of counts and identifies the additional statistic, the reflected difference set, needed for this operation. It does not give evidence about the performance of a memory architecture. It does give a finite example with exact counterexamples, a specified intervention, and an observable cost of discarding the statistic. That is the level at which this connection can be tested.
